\section{Two extracted transfer results}\label{sec:lemmas}

The following statements make two reusable parts of the proof explicit. They are proved here as mathematical deductions. Neither is claimed to be new in the literature, and neither constitutes a new resolution of a conjecture.

\subsection{A quantitative comparison certificate}\label{sec:certificate}

\begin{proposition}[Covolume comparison with explicit losses]\label{prop:certificate}
Let \(D>0\), \(m>0\), \(0<\alpha\le1\), and \(\alpha D\le n\le D\). Let \(M,U,L,I\ge0\), and let \(R,V,\delta\) be real numbers satisfying
\begin{equation}\label{eq:certificate-inputs}
D(M-L)\le V+2\delta,\qquad
V\le n(M-mR+U),\qquad
\delta\le DI.
\end{equation}
Then
\begin{equation}\label{eq:certificate-sharp}
nmR\le nU+D(L+2I)-(D-n)M,
\end{equation}
and in particular
\begin{equation}\label{eq:certificate}
R\le\frac{U}{m}+\frac{L+2I}{\alpha m}
\le\frac{U+L+2I}{\alpha m}.
\end{equation}
\end{proposition}

\begin{proof}
Combining the three hypotheses gives
\[
D(M-L)\le n(M-mR+U)+2DI.
\]
Rearrangement is exactly \eqref{eq:certificate-sharp}. Because \(D-n\ge0\) and \(M\ge0\), dropping the last term increases the right side. Divide by \(nm>0\), and use \(D/n\le1/\alpha\) together with \(L+2I\ge0\). The last inequality follows from \(U\ge0\) and \(\alpha\le1\).
\end{proof}

Here \(U\) is the loss in the upper norm estimate, \(L\) the loss in the lower estimate, and \(I\) the loss from denominators. This formulation identifies what a proposed adaptation must actually prove. Having a family of theta lifts is insufficient; its three normalized losses and its surviving rank must satisfy a quantitative certificate.

If, after all normalizations, the combined loss \(U+(L+2I)/\alpha\) is \(O_\eta(x^\eta)\) and \(m\ge x^\theta\), one obtains an upper bound of order \(x^{-\theta+\eta}\). This definition includes any loss from a varying spanning fraction \(\alpha\). With the particular Dirichlet quadratic-twist criterion from this repository, this excludes zeros only strictly to the right of
\[
\max\left\{\frac12,\,1-\theta+\eta\right\}.
\]
The \(1/2\) floor here belongs to that twist-selection theorem. No conclusion beyond its domain can be obtained just by increasing \(m\).

\subsection{Why the certificate is expensive to satisfy}

The apparently weak condition \(n\ge D/2\) is enough for the final division. It is not enough for the arithmetic lower bound that precedes it.

In \src{New/Determinant/Twist/LongForms.lean}, a missing basis vector has a completion cost bounded by
\[
2x_1^2b^2+b.
\]
The source supplies \(D-n\le4D/b^2\), yielding total cost at most
\[
4D(2x_1^2+b^{-1})\le Db
\]
under the chosen scale inequalities. A fixed positive missing fraction would instead leave a term of order \(Dx_1^2b^2\). Almost full spanning is therefore a mathematical necessity of this lower-bound argument, even though the last line of the proof uses only a positive fraction of the dimension.

The denominator comparison is equally informative. The allotted budget in \src{New/Determinant/Twist/FitTwists.lean} is
\[
\frac{fa(\log b)^3}{mk},
\qquad a=xb.
\]
The explicit local-defect estimate in \src{New/Congruences/FitTwistsExist.lean} has size
\[
O\!\left(f\frac{m}{k}b\log b\right).
\]
Their quotient has the form
\begin{equation}\label{eq:budget-ratio}
O\!\left(\frac{m^2}{x(\log b)^2}\right).
\end{equation}
Thus \(m\asymp\sqrt{x}\) leaves logarithmic room. Taking a larger power of \(x\) while keeping the subpower scale choices would require better loss estimates. More generally, if a normalized loss grows like \(m\), division by \(m\) leaves a constant error and the desired decay disappears.

\subsection{Only sufficiently many auxiliary primes are needed}\label{sec:subset}

The published criterion asks for the arithmetic bound at every permitted auxiliary prime. Its proof actually requires only one prime that is simultaneously arithmetically successful and analytically good for each possible zero.

\begin{theorem}[A large successful set suffices]\label{thm:subset}
Let \(q\) be a positive multiple of \(8\), let \(a\) be coprime to \(q\) with \(a\equiv7\pmod8\), and fix \(1/2<\beta<1\). Let \(S\) be a set of primes \(d\equiv a\pmod q\). Suppose:
\begin{enumerate}
\item For every \(d\in S\), \(R_{d,q}(x)\le x^{\beta-1}\) for all sufficiently large real \(x\). The threshold may depend on \(d\).
\item If \(A_S(X)=\#\{d\in S:d\le X\}\) and \(c=(\beta-1/2)/18\), then
\[
\frac{A_S(X)}{X^{1-c}}
\quad\text{is unbounded for positive integer \(X\).}
\]
\end{enumerate}
Then \(L(\rho,\chi)\ne0\) for every character modulo \(q\) and every \(\rho\ne1\) with \(\Rea\rho>\beta\).
\end{theorem}

\begin{proof}
Fix a proposed \(\rho\) and put \(T=|\Ima\rho|\) and \(\delta=(\beta-1/2)/9\). The theorem
\code{zero_density_quadTwist} in \src{Lib/Analytic/QuadraticTwist.lean} bounds by
\[
O_{q,\beta,T}\!\left(X^{1-\delta/2}\right)
=O_{q,\beta,T}(X^{1-c})
\]
the number of permitted primes \(d\le X\) for which some quadratic twist of a character modulo \(q\) has a zero in
\(\Rea s\ge\beta,\ |\Ima s|\le T\).

By the second hypothesis, \(S\) cannot be contained in this exceptional set. Choose \(d\in S\) outside it. For this fixed \(d\), the first hypothesis provides the required eventual one-sided bound. All the twisted multiplicities in \eqref{eq:residue} vanish in the chosen box. The remaining real residues are nonnegative; Landau's lemma therefore makes the residue at \(\rho\) zero. The denominator is nonzero and the numerator is a sum of nonnegative integers. In particular the multiplicity of a zero of the proposed \(L(s,\chi)\) must be zero, a contradiction.
\end{proof}

The order of quantifiers is essential: first choose \(\rho\), then a successful \(d\), then use the eventual estimate for that \(d\). No uniform bound on the thresholds \(X_d\) is required.

\begin{corollary}\label{cor:density}
Each of the following conditions on \(S\) is sufficient for the second hypothesis:
\begin{enumerate}
\item positive upper relative density among the primes in the permitted residue class;
\item \(\limsup_{X\to\infty}\log(1+A_S(X))/\log X=1\);
\item divergence of \(\sum_{d\in S}(\log d)/d\).
\end{enumerate}
\end{corollary}

\begin{proof}
For the first, the prime number theorem for a fixed progression gives counts on a sequence of order \(X/\log X\), and \(X^c/\log X\to\infty\). For the second, a bound \(A_S(X)=O(X^{1-c})\) contradicts the displayed upper exponent. For the third, such a counting bound would make the sum converge by partial summation, since the resulting tail integral is bounded by a constant times
\(\int_2^\infty (\log t)t^{-1-c}\,dt\).
\end{proof}

To recover GRH, use the theorem for every \(\beta\in(1/2,1)\); the successful set may depend on \(\beta\). A single set with counting exponent \(1\), on which the bound holds for every positive \(\varepsilon\) in the form \eqref{eq:one-sided}, also suffices.

This weakening could be useful when a new arithmetic construction works only for many auxiliary parameters rather than all of them. It is a reduction of required hypotheses, not a construction of those parameters or a new zero-free region by itself.
